% TOP_PROBLEM: {"id": 5300031, "title": "Components after removing a Cremer fixed point", "category_id": 11, "category": {"id": 11, "name": "dynamical_systems", "display_name": "Dynamical Systems", "description": "Problems about long-term behavior of deterministic systems, Hamiltonian dynamics, and periodic orbits.", "slug": "dynamical-systems", "order_index": 11, "created_at": "2026-07-31T15:26:25.671Z"}, "status": "partially_solved", "problem_number": "AMR-052-0031", "set_id": 13}
% FIRST_POSTED: 2026-10-02
% ENTRY_KIND: statement_only
% UnsolvedMath ID 5300031; source catalogue code AMR-052-0031
% !TeX program = lualatex
% Definition and sources only; this file is not an LLM solution attempt.
\documentclass[11pt,a4paper]{article}
\usepackage[margin=25mm]{geometry}
\usepackage{fontspec}
\setmainfont{lmroman10-regular.otf}[BoldFont=lmroman10-bold.otf,
 ItalicFont=lmroman10-italic.otf,BoldItalicFont=lmroman10-bolditalic.otf]
\usepackage{amsmath,amssymb,mathtools}
\usepackage{unicode-math}
\setmathfont{latinmodern-math.otf}
\AtBeginDocument{\let\setminus\smallsetminus}
\usepackage{xurl,hyperref}
\hypersetup{colorlinks=true,urlcolor=blue}
\setcounter{secnumdepth}{0}
\setlength{\parindent}{0pt}
\setlength{\parskip}{5pt}
\setlength{\emergencystretch}{3em}
\begin{document}
\section*{Components after removing a Cremer fixed point}\label{5300031}
{\small UnsolvedMath ID 5300031 · AMR-052-0031 · Dynamical Systems · AMR Open Problem Lists}
\subsection{Definitions and mathematical statement}
Let $P_\alpha(z)=z^2+e^{2\pi i\alpha}z$, where $\alpha$ is irrational
and the fixed point at zero is not locally holomorphically conjugate
to multiplication by $e^{2\pi i\alpha}$. Such a fixed point is called
a Cremer point. Let
$K_\alpha=\{z:\sup_{n\ge0}|P_\alpha^n(z)|<\infty\}$ and
$J_\alpha=\partial K_\alpha$ be the filled Julia set and Julia set.
The latter is considered with its ordinary subspace topology in
$\mathbb C$.

Remove just the fixed point and let
\[
 \mathcal C_\alpha=\{C:C\text{ is a maximal connected subset of }
                                  J_\alpha\setminus\{0\}\}.
\]
Determine the cardinality $|\mathcal C_\alpha|$ as $\alpha$ varies over
quadratic Cremer parameters. In particular, the source asks whether
this number is always countably infinite.

Connected components here mean maximal connected sets, not path
components, complementary domains in the plane, or components of a
small neighborhood of zero. These distinctions matter for a Julia set
which is not locally connected. Countably infinite means a bijection
with the positive integers; neither finite branching nor uncountably
many components satisfies that proposal.

A satisfactory answer may prove a universal value, give counterexamples
to it, or identify dynamical or arithmetic conditions under which the
number differs. The operation is deletion of one point from the entire
Julia set. Deleting the entire orbit of the critical point, taking a
quotient model, or replacing components by externally accessible
branches would ask different questions.
\subsection{Short English statement}
How many connected components remain when the Cremer fixed point is removed from a quadratic Cremer Julia set?
\subsection{Sources}
\begin{itemize}
\item[S1] ULAM.AI and the UnsolvedMath Contributors, supplied problems.json, record 5300031 (AMR-052-0031), AMR Open Problem Lists. Catalogue identity and source transcription; CC BY 4.0.
\url{https://huggingface.co/datasets/ulamai/UnsolvedMath}.
\item[S2] Stony Brook - Open Problems in Dynamical Systems, 1992, Milnor local P6. Original source indexed by AMR-052-0031.
\url{https://www.math.stonybrook.edu/open-problems-dynamical-systems}.
\item[S3] Milnor, Problems on local connectivity, Problem6, in Bielefeld–Lyubich (1992).
\url{https://www.math.stonybrook.edu/preprints/ims92-7.pdf}.
\end{itemize}
\end{document}
